\section{Evaluation framework}
\label{sec:evaluation}

This section describes the evaluation protocol. Each component was adopted in response to a
specific measured failure in an earlier cycle, and \S\ref{sec:background-vsix} describes those
failures.

\subsection{Outcome measures}
\label{sec:evaluation-measure}

The primary outcome is the team win rate, the proportion of games in which the evaluated
configuration's team takes at least five of the nine half-suits. Comparisons are reported as an
\textbf{edge in percentage points},
\[
  \text{edge} \;=\; 100 \times (\text{win rate} - 0.5),
\]
so that zero denotes a 50\% win rate in the specified comparison and the sign gives the direction.
The paper writes \emph{pp} for percentage points of this quantity, and a bracketed pair
$[\ell, u]$ after an estimate is its two-sided nominal 95\% confidence interval (CI) as
constructed in \S\ref{sec:evaluation-intervals}, which also states why every interval in this
paper is labelled nominal. A comparison ``against'' a named configuration is a head-to-head
comparison: an edge of $+3$~pp means a 53\% win rate against that configuration in the paired
design of \S\ref{sec:evaluation-duplicate}, not a three-point change from a separately measured
baseline win rate. Three other quantities appear and are never conflated with it: the mean number
of half-suits a team takes, used for the search's internal return; predictive log loss and argmax
hit rate, used to score belief models as predictors (\S\ref{sec:background-vsix}); and the
\emph{adversarial response win-rate edge} of \S\ref{sec:results-adversarial}, which is an edge in
the same units but measured for a fitted opponent against the evaluated configuration.

\subsection{Duplicate deal blocks}
\label{sec:evaluation-duplicate}

A \emph{cell} compares two configurations $A$ and $B$ over a \emph{bank} of deals identified by one
integer seed. Every deal in the bank is played twice: once with $A$ in the odd seats and $B$ in the
even seats, and once with the seats exchanged. Both configurations therefore play the identical deal
from both sides.

Fish has large per-deal variance, since a deal can be won or lost on the cards, and an unpaired
comparison spends most of its sample resolving that variance rather than the difference between
policies. The rotation balances team assignments within each deal and can reduce the variance associated
with deal strength, in the spirit of duplicate bridge and of the variance-reduction estimators
developed for poker \cite{zinkevich-divat,burch-aivat}. The design requires reproducibility from the seed alone
(\S\ref{sec:inference-engine}); without it the two rotations are not the same deal and the design
degrades silently into an unpaired one.

\subsection{Uncertainty estimation}
\label{sec:evaluation-intervals}

The two rotations of one deal are not independent observations: the two games share an initial deal
and are therefore treated as a single resampling block. Resampling blocks preserves their
within-deal dependence, whatever its sign.

Three interval constructions are used, and they are labelled by their construction rather than
described as one method.

\paragraph{Per-bank intervals.} Every cell measured on one bank carries a two-sided 95\%
  \textbf{deal-clustered percentile bootstrap interval}. The resampling unit is the \emph{deal},
  carrying both of its rotations together; 20{,}000 bootstrap resamples are drawn with replacement
  from the deals of the cell, the edge is recomputed on each resample, and the interval is the
  empirical 2.5th and 97.5th percentiles of that distribution. No Wilson or normal-approximation
  interval on the game count is quoted for a single bank. The harness records both intervals for
  every cell, and in this corpus the clustered interval is typically no wider than, and often
  slightly narrower than, the interval that treats games as independent, which is the direction a
  negative within-deal correlation produces.
\paragraph{Independence assumption.} All three constructions treat deal blocks as independent
  resampling units. In the scored implementation, agents are reused across the deals a worker
  thread processes and some rollout state persists between deals (\S\ref{sec:agent-coordination}),
  so a decision in one deal can depend on earlier deals and on thread scheduling, and the
  work-stealing schedule is not retained. For comparisons involving searching agents, then, the
  independence of deal blocks and the coverage of the resulting intervals have not been
  established; this is an unmodelled history and scheduling sensitivity, not a demonstration of
  miscoverage in any particular cell, and the non-searching cells are not affected by this
  mechanism. By convention every interval in this paper is labelled \emph{nominal}. The observed
  residual is small (\S\ref{sec:results-controls}), but there are only incidental observations of
  its effect, not a quantified bound, and the label does not repair it;
  \S\ref{sec:discussion-engineering} states what a corrected design would and would not establish.
\paragraph{Cross-bank pooling.} Two banks of equal size pool as the mean of their two edges, with
  standard error $\mathrm{se}_{\text{pooled}} = \sqrt{\mathrm{se}_1^2 + \mathrm{se}_2^2}/2$, where
  each $\mathrm{se}$ is recovered from that bank's bootstrap interval as half-width divided by
  $1.96$. The pooled interval is the estimate plus or minus $1.96\,\mathrm{se}_{\text{pooled}}$. This
  is a normal approximation built on the per-bank bootstrap standard errors, not itself a bootstrap
  interval. Both per-bank values are printed beside every pooled figure in this paper.
\paragraph{Between-cell contrasts.} A difference between two separately measured cells, such as an
  ablation's leave-one-out drop, is the difference of the estimates with the two half-widths
  combined in quadrature, omitting the covariance between the two cells. The configurations are
  paired on deals, so the two cells' estimation errors are generally correlated; the omission
  overstates the width only when that covariance is non-negative, which is expected when both cells
  are measured on the same deals against the same opponent but is not verified cell by cell. The
  harness produces no paired difference \emph{across} two separately executed cells, and the
  scored artifacts record per-cell summaries rather than per-deal outcomes, so a joint bootstrap of
  the actual contrast would require re-running the cells with per-deal capture; that is listed in
  \S\ref{sec:discussion-future}. Where a contrast combines cells that
  share banks and deals with unit coefficients, as the rule-variant residual of
  \S\ref{sec:results-dialects} does, the propagated interval is not a valid 95\% interval and is
  labelled as such where it appears.

\subsection{Sample-size planning}
\label{sec:evaluation-power}

At parity, the 95\% half-width of a cell's edge in percentage points is approximately
$98/\sqrt{N}$ for $N$ games, the half-width of a normal interval for an independent Bernoulli mean
at $p = 0.5$. The harness computes this planning approximation for every cell and emits it with the
result. It treats the $N$ games as independent and ignores the duplicate-block covariance, so it is a
planning figure rather than a power calculation: at $N = \vsevenScoreboardGames{}$ it gives a
one-point half-width, which does not mean that a one-point effect would be detected with any
specified probability. No formal power calculation, with a specified test, effect size, significance
level, desired power and block-level variance, was performed. The arithmetic is applied in advance
because it determines what may be attempted: a one-point half-width requires roughly
\vsevenScoreboardGames{} games, and a quarter-point half-width requires sixteen times that.

The failure this planning step addresses is the one \S\ref{sec:background-vsix} describes, in which
five mechanisms cleared a 95\% paired interval at one budget and returned exactly zero at three
times it. Sample sizes in this paper are chosen from the arithmetic before a cell is run, never
after seeing its result.

\subsection{Prespecified sign-agreement criterion}
\label{sec:evaluation-replication}

Every claim is measured on two disjoint deal banks, and the criterion was fixed before the data were
seen:

\begin{quote}
A claim whose \textbf{sign} does not agree on both banks is reported as not replicated, whatever its
pooled interval says.
\end{quote}

This is stricter than a pooled significance test, deliberately. A pooled interval can exclude zero
while the two banks disagree, and that configuration has occurred in this corpus and has previously
been published as a positive result. The criterion binds against this paper's own findings in
\S\ref{sec:results-attribution} and again in \S\ref{sec:results-panel}, and both instances are
reported.

\subsection{Holdout sealing and preregistration}
\label{sec:evaluation-seal}

A programme that generates candidate configurations, scores them and selects the best is performing a
maximisation, and the selected configuration's score is upward biased by the amount a maximum over
$K$ draws exceeds the mean. No care taken in the final measurement removes that bias if the final
measurement uses material the selection has seen.

Before any candidate existed, deal banks were reserved by seed and their digests committed, a digest
being a rolling hash over the six dealt hands and the dealer of every deal in the bank, computed
without constructing a policy or playing a game. An adversary population was split in half by a rule
fixed in advance, and one half was encrypted with the plaintext's SHA-256 committed, so that it could
not be trained against even accidentally. The engine's match command enforces the seal: a run
naming a sealed seed fails unless an explicit environment variable is set, which was set once, at
the start of the final evaluation, and recorded. The guard's scope is the match command; the
fitting, ablation and probe drivers derive their own seeds and bypass it, as Appendix~\ref{app:deviations}
(D6) records, and a committed digest is a commitment to content rather than a proof that the content
was not accessed. The documented access procedure is the one recorded there.

The entire final battery, comprising every cell, its sample size, its threshold, the arithmetic
behind the sample size and the conditions that would count as non-confirmation, was written down and
committed before any holdout bank had been played and before the sealed half had been decoded. The
evaluation was then executed as a separate automated agent session: a large language model
operating the repository under the author's supervision, with the session settings recorded in
\filepath{docs/v07/PHASE-PROMPTS.md}. The phase instructions required the session to start
without conversation history from the development phases and to read the registered protocol and
not the development-phase documents or training logs; the execution records document those
instructions and the session's outputs, not its compliance, which is not independently verifiable
from them. The restriction was by instruction rather than by access control, since the session had
read access to the repository, and the same author configured every phase; no evaluator
independence is claimed. Where the evaluation required a training figure for comparison it took it
from the registered protocol, which states such expectations in advance for exactly this purpose, so that agreement
constitutes a check and disagreement constitutes a finding.

Anything the evaluation did that the protocol does not specify is recorded as a deviation, next to
the affected result. There are \vsevenDeviations{}, listed in Appendix~\ref{app:deviations}.

\subsection{Selection effects and multiple comparisons}
\label{sec:evaluation-multiplicity}

This study evaluates many candidate mechanisms, opponents, ablations, adversarial searches and
configurations, and the inferential status of that needs stating precisely.

No formal family-wise or false-discovery correction was applied, and no interval reported
here is adjusted for multiplicity. What the design provides instead is structural.
\emph{Registration} fixes the primary comparison, its threshold and its interpretation before the
holdout outcomes were examined, so the primary claim is a single pre-specified test rather than the
best of many.
\emph{Sealing} guarded the named holdout banks against routine match access and committed their
digests before any candidate existed; the guard is an environment-variable override on the match
runner, some fitting and diagnostic drivers bypass it, and a digest commits to content without
making it inaccessible (\S\ref{sec:evaluation-seal}, Appendix~\ref{app:deviations} D5 and D6), so
this is a guarded procedure with prior commitments rather than unavailability by construction. \emph{Disjoint banks with a sign criterion} require a claim to appear twice
before it is reported. \emph{Calibration} (\S\ref{sec:evaluation-floor}) supports the interpretation of a null
adversarial result under the registered responder class, budget, handicap family and banks; it does
not make every null interpretable. And an explicit
\emph{selection check} (\S\ref{sec:results-selection}) computes how much of the measured gain a
maximum over $K$ candidates would produce under the null and compares advance-stated training figures
against their holdout counterparts.

Exploratory comparisons in \S\ref{sec:development} are reported as exploratory: they were run on
training material during the development phase, none is treated as confirmatory, and none is used to
support the paper's primary claim.

\subsection{Behavioural validity checks}
\label{sec:evaluation-gate}

Win rate and behavioural validity are dissociable in this policy family. The corpus contains a
configuration that scores \vsevenScoreboardTrapWin\% against an earlier version while making
\vsevenScoreboardTrapDead\% provably unsuccessful asks, asks whose target it can already prove does
not hold the card (\S\ref{sec:agent-policy}), and losing \vsevenScoreboardTrapLimit\% of its games
to the action cap. It improves win rate against the specified opponent while failing the chosen
behavioural criteria.

A configuration is therefore gated \emph{before} any strength number is computed for it, on
\vsevenGateRulesWord{} rules covering provably unsuccessful asks, runs of them, action-limit games,
the mirror game-length tail, late declarations, and the side-channel tests below. The protocol calls
this the \emph{soundness gate}; outside quotations this paper calls it the behavioural validity
gate, because it is an empirical test suite encoding behavioural criteria, not a proof. A provably unsuccessful ask is legal and
could in principle serve an information or turn-transfer purpose, so a rejection expresses a design
requirement rather than the invalidity of every rejected strategy. A configuration that fails is
rejected regardless of its win rate. The gate runs on \emph{training} material by design, being a
check on self-play behaviour rather than a strength claim. It is accompanied by a negative control,
a configuration constructed to be rejected, because a gate that cannot reject it would establish
nothing about the configuration it accepts.

\subsection{Adversarial-search calibration}
\label{sec:evaluation-floor}

An adversarial-response figure is produced by a search, and a search that finds nothing admits two
explanations: there is nothing to find within the class and budget, or the search is too weak.
Planted-handicap experiments characterise sensitivity to selected failure modes. An artificial
handicap of known size is planted in a target, a responder is fitted against the handicapped
target, and the \emph{excess} of that responder's edge over its edge against the unhandicapped
target is measured. The smallest planted edge the procedure recovers, for a given search class,
budget and handicap family, is that class's \textbf{detection floor}.

Three quantities that this number is sometimes made to stand for are kept separate in this paper.
The first is the observed planted-handicap recovery threshold for a particular search class,
budget and handicap family, which is what the calibration measures. The second is the sampling
uncertainty of a fixed policy comparison, which is a property of the cell's sample size and is
described by its interval (\S\ref{sec:evaluation-intervals}), not by the floor. The third is the
preregistered decision threshold of \vsevenPrimaryFloor{}~pp for declaring an advance
(\S\ref{sec:design-protocol}), which the protocol set equal to the phase-2 calibration floor of the
principal responder class. The protocol itself describes that use as a pre-committed reference bar
rather than a detection floor for a directly measured paired difference; it is a historical
decision rule, not a resolution limit for every comparison in the paper.

The \vsevenPrimaryFloor{}~pp figure has a specific origin. The phase-2 calibration planted a ladder
of handicaps of increasing size in a \vsix{} target and fitted a responder of the principal class at
a fixed budget against each; \vsevenPrimaryFloor{}~pp is the planted cost of the smallest rung
whose recovered excess excluded zero on both banks, and the next rung down, at
\vsevenFloorRungBelow{}~pp, was not recovered. These are two observed recovery outcomes under the
specified responder class, budget, handicap family and target; they do not estimate a continuous
detection threshold, and the registered \vsevenPrimaryFloor{}~pp decision threshold is a
historical choice.

Two findings from the calibration shape the adversarial statements in this paper. First, within the
tested class and budgets, the floor did not decrease with evaluation games: an early prediction that
it would fall as (evaluation games)$^{-1/2}$ was not borne out at its own predicted value, since at
four times the games the same rung, \vsevenPrimaryFloor{}~pp, was the smallest recovered and the
next rung down remained unrecovered, with a phase-2 excess whose interval was centred near zero
rather than merely wide. This is one comparison, of two evaluation sizes for one responder class on
nested banks. For the tested responder class at these
budgets, then, more evaluation games and more fitting computation address different sources of
uncertainty. The exact result is that the \vsevenFloorRungBelow{}~pp rung was not recovered whereas
the \vsevenPrimaryFloor{}~pp rung was; no continuous detection threshold between them is estimated.
Failure to recover a small planted effect does not establish its absence, and recovery of one
handicap family does not establish sensitivity to every exploit of the same size. Second, floors are
family-specific: the declaration channel has its own and higher floor of \vsevenDeclFloor{}~pp,
measured as a win-rate response to a declaration handicap, and reading a declaration-objective
result against the general floor would overstate it. That figure is in win-rate units and is not
applied to declaration accuracy without a justified mapping of units.

An adversary class contributes to an adversarial-response statement only after passing a
planted-edge calibration at or below the effect size in question. That precondition is why the
negative controls of \S\ref{sec:results-controls} are run: an instrument that cannot recover a
planted edge of a given size cannot report the absence of a real one of that size.

\subsection{Information-access and isolation tests}
\label{sec:evaluation-side}

The evaluated team is three identical agents with no shared secret. That their policy specification
is common knowledge is legitimate (\S\ref{sec:game-why}), but it also means an implementation detail
could become an unintended signalling channel. Before this cycle the constraint had only ever been
checked by reading the source.

Four mechanical tests now run over all four engine decision types. \textbf{S3, listening
substitution} asks whether a seat's behaviour changes more when a teammate's hidden action changes
than when an opponent's does, beyond what the rules explain. \textbf{S4, stream independence} asks
whether the seat's private random stream is drawn independently of the deal. \textbf{S5, posterior
invariance} asks whether the seat's action depends on true card locations beyond what its own
posterior supports. \textbf{S6, seat isolation} replays every decision from own hand, public event
stream, rules and reset seed alone, and counts the decisions the reconstruction cannot reproduce, at
zero tolerance.

The tests are calibrated by three positive controls, each a policy carrying one planted violation,
and each must fail exactly the tests the protocol names for it in advance and no others.
\S\ref{sec:design-material} reports the calibration on holdout material. Without those controls a
clean report from the gate would be uninformative.

One limitation is stated wherever these tests are used. The zero-tolerance form of S6 requires a
single worker thread, because the audit depends on the deal-to-thread partition and work stealing
does not fix that partition at higher counts, and it requires agents rebuilt for each deal.
Rebuilding agents per deal \emph{changes play}. The fresh-agent, single-thread configuration is
therefore the one that passed the specified isolation tests. The scored configuration, with reused
agents and \vsevenScoredThreads{} worker threads, retained a cross-deal state residual in searching
agents,
so these test results do not establish strict seat isolation or exact reproducibility for the scored
mode, and every strength number in this paper is measured under that mode
(\S\ref{sec:results-controls}).

\subsection{Worst-case performance and empirical regret}
\label{sec:evaluation-panel}

The project's standing reporting rule is that the worst case leads and an aggregate never does, since
a mean over an opponent panel rewards beating weak opponents by larger margins, which is not the
quantity of interest. Two statistics are reported over a panel scored identically for every
configuration. Write $e(a,o)$ for the pooled edge of configuration $a$ against panel member $o$,
over the four evaluated configurations $a$ and the fixed panel $\mathcal{O}$. The \textbf{worst
cell} of $a$ is $\min_{o \in \mathcal{O}} e(a,o)$, and its \textbf{empirical minimax regret} is
\[
  R(a) \;=\; \max_{o \in \mathcal{O}} \Bigl[\, \max_{b} e(b,o) \;-\; e(a,o) \Bigr],
\]
the largest shortfall from the best of the four configurations on any member. Both are defined over
the evaluated configurations and the fixed panel; neither is a global minimax guarantee. They can
disagree, and in \S\ref{sec:results-panel} they do. We report both statistics because they assess
different aspects of panel performance, and report the panel mean last and labelled as a
diagnostic. Extrema and rankings computed from point estimates are descriptive: the interval printed
beside a selected worst cell is that cell's own interval, not a confidence interval for the minimum
over the panel, and no simultaneous or joint-bootstrap inference over the panel is supplied.

Every configuration must carry a cell for every member or the regret is undefined, which is why the
panel is fixed by name in advance and scored for all configurations on the same banks. The
\vsevenPanelMembers{} named members are \vsevenPanelDistinct{} distinct policies, since three
sealed aliases duplicate members that also appear under their own names. Three of the evaluated
configurations are themselves panel members and meet themselves; a self-cell is a mirror match with
an edge of exactly zero and no sampling variance. Under the higher-is-better convention a zero cell
can only lower a configuration's worst cell, which it does exactly when every other cell is
positive, and it can only raise a configuration's regret, so \S\ref{sec:results-panel} prints the
worst case both with and without self-cells so that the two can be compared; the reported regrets
are unaffected because none is attained at a self-cell.
